% M4 PAPER-NUCLEO --- the axiomatic core of the programme (ROADMAP_V2 item M4).
% v0 born from N0+N0' with the M1 (combinatorial trichotomy) and M2/N5 (radiative
% Skyrme) sections already hardened.  Claim discipline (charter): every statement
% is about the CLASS defined by the axioms, never "Nature obeys"; [measured]
% inputs are cited as measurements of the programme; open flanks are a section,
% not a footnote.  Compiles with pdflatex+bibtex (revtex4-2).
\documentclass[aps,prd,reprint,nofootinbib,superscriptaddress,floatfix]{revtex4-2}
\usepackage{amsmath,amssymb,amsthm}
\usepackage{bm}
\usepackage{booktabs}
\usepackage{hyperref}

\newtheorem{theorem}{Theorem}
\newtheorem{lemma}{Lemma}
\newtheorem{proposition}{Proposition}
\newtheorem{corollary}{Corollary}
\theoremstyle{definition}
\newtheorem{axiom}{Axiom}
\newtheorem{definition}{Definition}
\newtheorem{prediction}{P\!}

\begin{document}

\title{What a Lorentz-invariant discrete order can carry:\\
an axiomatic core for emergence on causal substrates}

\author{Miqueias Alves Mendes}
\email{mikalvesm@gmail.com}
\affiliation{Independent Researcher, Ibiapina, Cear\'a, Brazil}
\date{\today}

\begin{abstract}
We organize a research programme's results into an axiomatic core. Two axioms---a
Poisson-sprinkled causal order (forced, not chosen: it is the unique Lorentz-invariant
discretization) and a generic compact internal field---define an \emph{intrinsic class}
of theories. One lemma fixes what a pair of events carries: time with sign and structure,
space without sign and without structure. One theorem then reads the isometry group
through three structural layers: non-compact orbits force divergent valence (no scale,
no emergent photon); disconnected components cannot condense (no parity violation, no
spontaneous arrow); the compact part protects discrete indices (topological charge,
exact multiplets, $\mathbb{Z}_2$ exchange statistics). The loop side of the argument is
upgraded from measurement to a combinatorial trichotomy: exchangeable orders have no
locally finite Hasse diagram; growth orders confine all cycles to finite blocks between
posts; and square-generated percolating loop structure \emph{forces} a grading---an
intrinsic foliation---so it cannot coexist with boost invariance. Renormalization is
built in: Poisson is the exact fixed point and global attractor of thinning, with no
invariant relevant direction; the density is the theory's single external unit. The
theory keeps exactly two external inputs beyond the axioms, each with a proof of
necessity: the Skyrme-type core cost, whose dominance cannot emerge at tree level (a
pointwise identity) and does not emerge radiatively (a measured one-loop channel
coefficient, negative at $\sim$130$\sigma$); and the lattice-to-SI conversion, one number
because there is one unit. We list the falsifiable predictions the core generates, the
places where it demonstrably does not deliver, and the ten open frontiers---including
the single combinatorial flank (rank-frustrated loop lattices) that our theorems do not
close.
\end{abstract}

\maketitle

\section{Introduction: the inverted question}
\label{sec:intro}

A research programme that couples matter fields to a causal set eventually accumulates
two kinds of results: positives (long-range order, topological solitons, an exact meson
octet, confinement, an emergent Newtonian sector~\cite{MatterGravityPRD,SU3PRD}) and
negatives (no emergent photon, no scale, no chirality, no spontaneous arrow of time).
For most of our programme's history the question was \emph{does $X$ emerge?}, answered
one campaign at a time. This paper inverts the question: \emph{what is a
Lorentz-invariant discrete order permitted to carry?} The answer turns out to be
compact---in both senses of the word---and it converts the programme's scattered record
into a short deductive core: two axioms, one lemma, one boundary theorem with three
layers, a combinatorial trichotomy, a built-in renormalization statement, and exactly
two external inputs each carrying a proof of its own necessity.

Three disciplines govern every claim below. First, all statements are about the
\emph{class} defined by the axioms---never about Nature; contact with Nature is listed
separately as predictions (Sec.~\ref{sec:predictions}). Second, statements carry their
epistemic grade explicitly: [theorem] (proved here or cited), [sketch] (complete
argument, gaps declared in Sec.~\ref{sec:open}), or [measured] (a pre-registered
campaign of the programme). Third, the places where the principle does \emph{not}
deliver are a section of this paper (Sec.~\ref{sec:vaccine}), because an exclusion
principle that explained everything would explain nothing.

A companion paper~\cite{WenComplement} develops one corollary---why string-net
condensation does not transport to causal substrates---at full length; here it appears
as one row of a table.

\section{Axioms and the intrinsic class}
\label{sec:axioms}

\begin{axiom}[substrate]
The substrate is a causal set $(C,\prec)$ obtained by a Poisson sprinkling of density
$\rho$ into $\mathbb{M}^d$ ($d\ge2$), with the discrete--continuum correspondence
frame-independent.
\end{axiom}

Poisson is not an independent postulate: it is the unique Lorentz-invariant discrete
measure (Bombelli--Henson--Sorkin~\cite{BHS2009}); the contrast is also measured in the
programme (a regular lattice fails isotropy) [measured]. Discreteness plus invariance
therefore \emph{forces} Axiom~1 up to the single parameter $\rho$.

\begin{axiom}[matter]
On $C$ lives a generic internal field $n\colon C\to X$ with $X$ compact (a group
manifold $G$ or coset $G/H$, $G$ compact), coupled by an invariant functional of the
intrinsic data (order and counting); genericity fixes the leading term
$\propto(1-n_i\!\cdot n_j)$ on related pairs [literature: universality; robustness to
$\pm10\%$ deformations measured].
\end{axiom}

\begin{definition}[intrinsic class]
An \emph{intrinsic} rule or observable is a measurable functional of $(C,\prec,n)$
alone---order, counting, and the internal field; no embedding coordinate.
\end{definition}

\begin{lemma}[pair content]
\label{lem:pair}
{\rm[theorem; proof in Appendix~\ref{app:proofs}]} For two sprinkled events: if $i\prec j$ (timelike), the intrinsic datum is
\emph{signed and structured}---the relation has a direction, and the Alexandrov interval
$[i,j]$ carries internal structure: $N(i,j)\sim\mathrm{Poisson}(\rho c_d\Delta\tau^d)$
and the longest chain $\approx m_d\rho^{1/d}\Delta\tau$
(Brightwell--Gregory~\cite{BrightwellGregory1991}; the limit exists by Kingman's
subadditive theorem) estimate proper time \emph{with sign}. If $i\parallel j$
(spacelike), the intrinsic datum is \emph{one bit} (incomparability); all further
information is third-party reconstruction and yields only \emph{unsigned} invariants
(estimates of $|s^2|$~\cite{RideoutWallden2009}, Gram data), determining a spatial
configuration at best up to the full orthogonal group $O(d-1)$, reflections included.
\end{lemma}

\noindent\textbf{Slogan.} \emph{The causal order carries time with sign and space
without sign---and without structure.}

Two anchors calibrate the lemma. In the continuum, Malament and
Hawking--King--McCarthy~\cite{Malament1977,HKM1976} show that causal order determines
the metric up to a conformal factor, fixed by volume: ``order + number = geometry''.
Lemma~\ref{lem:pair} is the finite-density form of that theorem: the timelike half is
intrinsic at any density; the spacelike half is asymptotic-statistical, recovered only
in the infinite-information limit. (The full discrete recovery is the causal-set
Hauptvermutung---a conjecture, with instances measured in the programme: dimension,
curvature at $23.5\sigma$, Schwarzschild proper time to $0.2\%$ [measured].) Second, an
antichain---the discrete support of a Cauchy surface---is \emph{amorphous}: its induced
suborder is empty, its automorphism group is the full symmetric group, and its only
intrinsic datum is its cardinality [theorem, trivial but central]. Spacelike $p$-cells
with intrinsic structure and orientation do not exist in the class, for any $p\ge1$.

\section{The Boundary Theorem}
\label{sec:boundary}

\begin{theorem}[boundary]
\label{thm:boundary}
{\rm[assembled; component proofs in Appendix~\ref{app:proofs} and citations]} In the
intrinsic class (Axioms 1--2), emergent
statistics sees the group $O(d-1,1)\times G_{\rm int}$ through exactly three structural
layers:
\begin{enumerate}
\item \textbf{Non-compact orbits $\Rightarrow$ divergence.} Every invariant of a causal
pair is a function of $\Delta\tau$ (Lemma~\ref{lem:pair}); the orbit of ``neighbours at
fixed $\Delta\tau$'' is the hyperboloid $H^{d-1}=SO^{+}(d-1,1)/SO(d-1)$, of infinite
volume; by Campbell--Mecke,
$\langle z\rangle=\rho\,\mathrm{Vol}(H^{d-1})\int\Delta\tau^{d-1}q(\Delta\tau)\,
d\Delta\tau=\infty$ for every non-trivial invariant rule [theorem]. Refused: finite
valence, correlation length, genuine criticality, emergent scale.
\item \textbf{Disconnected components $\Rightarrow$ unfixable bits.}
$\pi_0(O(d-1,1))=\mathbb{Z}_2\times\mathbb{Z}_2$ (parity $\times$ time reversal). The
$P$ component acts \emph{trivially} on intrinsic data (it is an isomorphism of
$(C,\prec)$); the $T$ component acts as \emph{order duality} ($\prec\,\leftrightarrow\,
\succ$), statistically symmetric. Refused: spatial orientation (even spontaneously,
Sec.~\ref{sec:parity}), arrow of time (statistically; the spontaneous channel is closed
by measurement).
\item \textbf{Compact part $\Rightarrow$ protected discrete indices.} Invariants of
compact groups have discrete spectra: $\pi_3(G)=\mathbb{Z}$ for every simple compact $G$
(Bott), $\pi_4(SU(2))=\mathbb{Z}_2$, representations of compact $H$ discrete and
finite-dimensional. Permitted---and, when the dynamics orders, protected: topological
charge, exact multiplets, $\mathbb{Z}_2$ exchange statistics, Goldstone counting.
\end{enumerate}
\end{theorem}

\begin{table}[b]
\caption{\label{tab:corollaries}The named corollaries and what each uses. All four
negative corollaries were first found as measurements; the theorem is what makes them
one fact.}
\begin{ruledtabular}
\begin{tabular}{llp{3.1cm}}
Corollary & Uses & Extra input \\
\colrule
no-scale & layer 1 & none (clean) \\
no-photon & layers 1+2 & $B\equiv$ spacelike flux (definition) \\
no-chirality & layer 2 ($P$) & none (Sec.~\ref{sec:parity}) \\
axis-without-arrow & layer 2 ($T$) & spontaneous side: [measured] \\
topological matter & layer 3 & Axiom 2; existence [measured] \\
\end{tabular}
\end{ruledtabular}
\end{table}

The scope is declared, not hidden: the principle is \emph{kinematic}. It governs what
the class can express and the \emph{form} of what emerges (which indices, which
degeneracies); the \emph{existence} of the positives (long-range order at finite $J$,
confinement) is dynamical and was measured, not deduced. An exclusion principle should
not predict phase diagrams.

\emph{The Wigner precedent.} The pattern is already standard physics: in Wigner's
classification~\cite{Wigner1939}, numbers that lock discrete (spin, helicity) are
invariants of compact little groups, and labels that run continuous (momentum on the
mass shell $H^3$) live on non-compact orbits---massless continuous-spin representations,
never observed, are exactly the non-compact part acting non-trivially. The same
hyperboloid that carries Wigner's continuous labels kills our valence. The Boundary
Theorem is the statistical version, over order, of a pattern quantum field theory
already practices.

\section{The combinatorial trichotomy: loops are theorems now}
\label{sec:trichotomy}

Layer~1 forbids finite valence for \emph{invariant rules on the sprinkling}. A
string-net substrate~\cite{LevinWen2005,WenComplement} needs more: finite valence
\emph{and} a percolating, finite-dimensional space of closed loops. Our programme
measured seven substrate families fail one or the other, never neither---the ``binary
alternative''. The following results replace that measured binary, in the three named
invariance classes, by theorems. Fix a poset $(P,\prec)$ with Hasse diagram $H$ (covering
relations $x\lessdot y$), and for a closed walk $\gamma$ let the \emph{circulation}
$\varphi(\gamma)$ be (up-steps $-$ down-steps).

\begin{lemma}
\label{lem:balanced}
{\rm[theorem]} Hasse diagrams have no triangles; every 4-cycle of $H$ has $\varphi=0$
(the directed square is acyclicity-forbidden; the $\varphi=\pm2$ pattern contradicts
covering).
\end{lemma}

\begin{theorem}[hinge: square-generated loops force a foliation]
\label{thm:T1}
{\rm[theorem]} Let $G\subseteq H$ be connected with cycle space over $\mathbb{Z}$
generated by 4-cycles. Then there is $r\colon V(G)\to\mathbb{Z}$, unique up to a
constant, with $r(y)=r(x)+1$ on every covering edge of $G$: the suborder is graded.
\end{theorem}

\begin{proof}
By Lemma~\ref{lem:balanced}, $\varphi$ annihilates the generators, hence the whole cycle
space over $\mathbb{Z}$; the path integral $r(x)=\varphi(x_0\!\to\!x)$ is
path-independent and rises by 1 on each up-edge.
\end{proof}

The hypothesis is exactly Wen's requirement, with no slack: Wilson loops decompose
multiplicatively into \emph{oriented} plaquettes---integer generation, not mod-2. The
level sets $\{r=k\}$ are antichains for the order generated inside $G$, and $r$ is
built from order alone, so every automorphism preserving $G$ preserves $r$ up to a
constant: \emph{a macroscopic square-generated cluster carries an intrinsic,
equivariant, discrete foliation---a preferred time readable off the order}. This is the
hinge of the programme's ``seven deaths'' made exact: arming square plaquettes
\emph{obliges} the foliation that breaks boost invariance; the CDT-like family sits on
that side by construction, and nothing invariant can sit on it at all.

\begin{theorem}[exchangeable exclusion]
\label{thm:T2}
{\rm[theorem]} No exchangeable random order on $\mathbb{N}$ has a non-empty, locally
finite Hasse diagram: by Aldous--Hoover/Janson~\cite{Janson2011}, the covering array is
exchangeable with constant density $p$ per ergodic component; $p>0$ gives every element
infinite expected Hasse degree, $p=0$ gives an a.s.\ empty covering relation.
\end{theorem}

\begin{proposition}[box orders]
\label{prop:T3}
{\rm[theorem, cited anchor]} For the dominance order on $n$ i.i.d.\ points in $[0,1]^d$ (for $d=2$,
the causal order of $\mathbb{M}^2$ at $45^\circ$), the expected number of covering pairs
is $\sim\frac{2}{(d-1)!}\,n(\log n)^{d-1}$, so
$\langle z\rangle_{\rm Hasse}=\Theta((\log n)^{d-1})\to\infty$~\cite{Brightwell1993}.
Even without the Lorentz group, order in $d\ge2$ makes the divergence inevitable in weak
(logarithmic) form; the boost upgrades it to a power.
\end{proposition}

\begin{theorem}[growth orders: cycles confined by posts]
\label{thm:T4}
{\rm[theorem, given cited anchors]} In transitive percolation (the simplest
Rideout--Sorkin sequential growth model~\cite{RideoutSorkin2000}), posts occur with
positive density and the infinite order decomposes into an ordinal sum of i.i.d.\ finite
blocks~\cite{ABBJ1994,BollobasBrightwell1997}. No Hasse edge crosses a post (if
$x\prec w\prec y$ strictly, $x\lessdot y$ is impossible) [theorem, new but elementary],
so the Hasse diagram is a chain of finite blocks joined at cut vertices: every cycle is
confined to a single $O_P(1)$ block, the number of independent cycles in radius $R$
grows $\lesssim R$, and loop structure never percolates in dimension $\ge2$.
\end{theorem}

\noindent\textbf{Trichotomy.} A random order carrying percolating finite-dimensional
loop structure would need locally finite valence (against Theorems~\ref{thm:T2} and
Proposition~\ref{prop:T3} in the exchangeable/embedded classes), unconfined cycles
(against Theorem~\ref{thm:T4} in the growth class with posts), and---if the plaquettes
are squares---no foliation (against Theorem~\ref{thm:T1}, which forces one). The
intersection demanded by string-net emergence is empty in all three named invariance
classes.

\noindent\textbf{Two retrodictions [measured].} (i) Sprinkled orders are violently
rank-frustrated (chains of different lengths between the same endpoints at every scale),
so by the contrapositive of Theorem~\ref{thm:T1} square-generated clusters cannot grow:
measured as flat, non-growing ``diamond towers'' ($\sim0.25\%$). (ii) In sequential
growth, intra-block cycles exist but never percolate: measured as the sub-mean-field
$C_4$ plateau ($\approx0.019$ at $N=16000$) of the growth-model campaign.

\noindent\textbf{The declared flank.} Theorem~\ref{thm:T1} governs balanced (square)
generators only. A homogeneous, finite-valence order whose cycle space is generated by
\emph{rank-frustrated} cycles (pentagons, $\varphi=\pm1$) and percolates in two
dimensions of cycles would escape all three legs. We can neither construct nor forbid
it; it is the programme's named open target (Sec.~\ref{sec:open}, item 1). Two
observations narrow it: every foliated/CDT-family substrate uses squares, so a
pentagonal substrate would be new physics by itself; and percolating rank frustration is
precisely the profile of sprinklings, where valence diverges---suggesting, not proving,
that frustration and finite valence compete.

\section{The signature boundary}
\label{sec:signature}

Where exactly does the refusal begin? Not at ``Lorentzian vs.\ Euclidean'' but at the
\emph{definiteness of the signature} [sketch]: for an invariant point process on
$(\mathbb{R}^d,g)$ with $g$ of signature $(p,q)$ and any invariant non-trivial connection
rule, the orbit of neighbours at fixed invariant separation is the quadric
$\{x\!\cdot\!x=c\}$. For definite signature the orbit is the compact sphere $S^{d-1}$
and $\langle z\rangle<\infty$ for every integrable profile---divergence is
\emph{chosen}; for indefinite signature every separation quadric has infinite volume and
$\langle z\rangle=\infty$ for every rule---divergence is \emph{forced}. The Lorentzian
case $q=1$ is merely the unique indefinite signature with a convex cone, i.e., with an
\emph{order}. Wick rotation is exactly the operation that crosses the boundary (it
compactifies the group), and an order cannot cross it and remain an order. This is why
sixty years of condensed-matter emergence---which lives entirely on the definite
side---do not transport to causal substrates; the companion paper~\cite{WenComplement}
develops the consequences.

\section{Renormalization is built in: the thinning fixed point}
\label{sec:thinning}

The coarse-graining native to the class is random deletion (thinning) $R_p$: keep each
element independently with probability $p$, take the induced suborder~\cite{Surya2019}.
In $\mathbb{M}^d$ the dilation $x\to x/b$ is an automorphism of the causal structure and
maps $\mathrm{PPP}(\rho)$ to $\mathrm{PPP}(\rho b^d)$, so \emph{the law of the intrinsic
order is $\rho$-independent}: $\rho$ is the theory's external unit, and the rescaling
step of a usual RG acts trivially. Three statements then close the scale question from
inside:

\begin{proposition}[exact fixed point]
{\rm[theorem]} $R_p(\mathrm{PPP}(\rho))=\mathrm{PPP}(p\rho)$; at the level of the
intrinsic order law, the sprinkling is an \emph{exact} fixed point of deletion---and by
BHS uniqueness, the only Lorentz-invariant one.
\end{proposition}

\begin{proposition}[global attractor, at the level of laws]
{\rm[theorem]} By the thinning-limit theorems~\cite{Kallenberg2017}, iterated
compensated thinning converges \emph{in distribution} to a Cox process directed by the
limiting sample intensity; for an \emph{ergodic} invariant process of finite intensity
$\lambda$ the spatial ergodic theorem makes the directing measure the constant
$\lambda$, and $\mathrm{Cox}(\lambda)=\mathrm{PPP}(\lambda)$. The fixed point is not
only unique: it attracts the whole class, in law. The finite-intensity qualifier is
load-bearing but harmless: an invariant cluster dressing would need an invariant
probability measure on a non-compact Lorentz orbit, and the only invariant probability
on $\mathbb{M}^d$ is $\delta_0$ (Weil uniqueness; Appendix~\ref{app:proofs})---the same
non-compactness of layer~1 protects the attractor.
\end{proposition}

\begin{proposition}[no relevant direction, with eigenvalues]
{\rm[theorem; proof in Appendix~\ref{app:proofs}]} The compensated step acts on
factorial moment densities as \emph{pure dilation of arguments},
$\rho_k\mapsto\rho_k(\cdot/a)$ with $a=p^{1/d}$: homogeneous truncated-correlation
directions $g\sim|x|^{-\alpha}$ are eigenfunctions with eigenvalue
$p^{\alpha/d}<1$ for every $\alpha>0$---all decaying tails are irrelevant with explicit
rate; the constant direction ($\alpha=0$) is intensity superselection (a Cox mixture),
intrinsically invisible in $\mathbb{M}^d$; the only marginal direction is $\rho$ itself,
the unit. Growing correlations are excluded by existence, not classified [declared].
\end{proposition}

Emergence is what survives deletion (topology, forms---discrete invariants cannot
flow); non-emergence is what deletion dilates away (correlations, scales, loops). The
one door the programme had left open for a new scale mechanism---a relevant direction
under thinning---is closed from inside the RG.

\section{Discrete spacetime symmetries: the parity no-go and the weaker $T$ side}
\label{sec:parity}

\begin{lemma}[intrinsic action of isometries]
\label{lem:isometries}
{\rm[theorem]} Every isometry of $\mathbb{M}^d$ preserving time orientation (rotations,
boosts, translations, \emph{and spatial reflections}) preserves $\prec$ and hence acts
\emph{trivially} on intrinsic configurations: it is descriptive redundancy, not a
symmetry of the intrinsic theory. Time-reversing isometries act as the \emph{order
duality} $(C,\prec)\mapsto(C,\prec^{\rm op})$, generically a distinct intrinsic object.
\end{lemma}

\begin{lemma}[spatial orientation has no carrier: exhaustion]
\label{lem:gram}
{\rm[theorem; proof in Appendix~\ref{app:proofs}]} Any intrinsic observable is a
measurable function of the isomorphism class of $(C,\prec,n)$; every
time-orientation-preserving isometry---spatial reflections included---induces the
\emph{identity} on isomorphism classes, so any putative $P$-odd intrinsic functional
satisfies $O=-O$ and vanishes identically. No enumeration of reconstruction schemes is
needed. Precisely: no \emph{absolute} sign exists; the \emph{relative} chirality of two
rigid clusters is intrinsic but $P$-\emph{even} (the product of two signs survives the
one global reflection).
\end{lemma}

\begin{lemma}[no chiral long-range order]
\label{lem:quenched}
{\rm[theorem for local estimators; proof in Appendix~\ref{app:proofs}]} For any two
disjoint bounded regions and any local chirality estimators $\chi_A,\chi_B$ (functionals
of the points in each region, odd under a time-preserving spatial reflection of that
region), $\mathbb{E}[\chi_A\chi_B]=\mathbb{E}[\chi_A]\mathbb{E}[\chi_B]=0$ at
\emph{every} separation: complete independence of the Poisson restriction plus local
reflectability---exclusion by \emph{independence of the measure}, not by symmetry.
Declared residue: estimators entangled with a common frame cluster are not covered
[sketch]; for dynamical geometry, outside Axiom~1, the lemma does not apply---declared
flank.
\end{lemma}

\begin{theorem}[parity no-go, strong form]
\label{thm:parity}
In the intrinsic class: (i) explicit $P$ violation is \emph{inexpressible} (no odd
functional exists to write); (ii) spontaneous $P$ violation is \emph{impossible}:
breaking requires an order parameter on which the symmetry acts, and $P$ acts trivially
on the whole intrinsic configuration space (Lemma~\ref{lem:isometries})---a redundancy
cannot break, in the same spirit as Elitzur's theorem~\cite{Elitzur1975}; (iii) declaring
the internal field pseudo-scalar is operationally empty: observable $P$ violation would
be an odd correlator between internal and \emph{geometric} handedness, and the geometric
partner has no carrier (Lemmas~\ref{lem:gram},~\ref{lem:quenched}). Hence emergent
chirality---correlation of internal with spatial handedness---does not emerge from the
class, explicitly or spontaneously.
\end{theorem}

The $P$-vs-$T$ structure is \emph{asymmetric}, and the theorem must say so. $P$ has no
carrier: the exclusion is kinematic and closes both channels. $T$ \emph{has} a carrier
(order duality is intrinsic; $T$-odd functionals exist), the law is self-dual so all
$T$-odd expectations vanish in statistical equilibrium, but spontaneous $T$ breaking is
not kinematically excluded---it was \emph{measured absent}: the axis (comparability)
emerges; the arrow does not condense and must be input [measured]. The $T$ closure is
empirical, and is declared as such. Finally, define $C$ as an internal involution that
is a symmetry of the action: $C$ acts non-trivially and \emph{may} break spontaneously
(layer 3); since $P$ acts trivially, $CP\equiv C$ operationally. Any $CP$ violation
whose content is genuinely $P$-entangled is therefore external input in this class.

Three sectors of the Standard Model are precisely the three chiral ones: neutrinos,
the generation structure, the electroweak gauge sector. In this class they are not
missing campaigns; they are missing \emph{carriers}. And the no-go is attackable:
any derivation of chirality from an invariant causal substrate must break
Lemma~\ref{lem:isometries} or Lemma~\ref{lem:gram}.

\section{The matter sector: what compactness protects}
\label{sec:matter}

Layer 3 is where the programme's positives live, and the theorem fixes their form. There
is no single universal index: there is a \emph{family}---$\pi_3(G)=\mathbb{Z}$ (baryon
charge; Bott for every simple compact $G$), $\pi_4(SU(2))=\mathbb{Z}_2$
(Finkelstein--Rubinstein spin-statistics), representations of unbroken $H$ (the exact
octet), $\dim G/H$ (Goldstone counting)---with one common root: \emph{compactness makes
spectra discrete}. All four were measured in the programme
[measured,~\cite{MatterGravityPRD,SU3PRD}], and the theorem says why the list looks like
that and not otherwise.

Two scope statements keep the section honest. First, the compact target is
\emph{hosted, not selected}: replacing $SU(3)$ by $SU(4)$ reproduces the full
phenomenology (ordering at finite $J$, $15/15$ gapless Goldstones, integer $B$,
confinement with decreasing $\sigma(\beta)$) [measured]---the ``3'' is the world's
choice, not the substrate's; whether confinement depends on the \emph{center} $Z(G)$ is
an open, discriminating test ($G_2$, trivial center, queued). Second, existence is
dynamical: the ordered vacuum at finite coupling is a measurement, not a corollary.

This division now survives quantization. Applying the forced covariant (Sorkin--Johnston)
prescription (Sec.~\ref{sec:predictions}, P1) to the orientation field over the
\emph{measured} ordered vacuum preserves the Goldstone multiplet: the quantum
level-splittings track the classical ones (ratios $\approx1.1$--$1.4$ against a $3\times$
anomaly threshold, robust to a fourfold change of mode count) and the relativistic branch
degrades on a disordered background, so the ordered vacuum is the carrier [measured,
pre-registered]. Crucially the SJ spectrum stays non-degenerate (minimum gap
$\sim\!10^{-3}$), so the emergent stabilizer is abelian $U(1)^K$ and \emph{no compact fiber
is created}: the quantization respects the postulated $G$ without generating it. Axiom~2
therefore remains a postulate---a delimitation, now measured rather than assumed---and the
rule ``forms are derived, groups are postulated'' holds intact into the quantum sector.

\section{The two external inputs, each with a proof of necessity}
\label{sec:externals}

The frozen theory is: \emph{two axioms plus two external inputs with proof}. ``With
proof'' means: for each input we can demonstrate the class does not supply it---so
importing it is not an omission but a theorem-backed boundary of the theory.

\subsection{The Skyrme-type core cost}
\label{sec:skyrme}

The Skyrme \emph{operator} is not external: the fourth order of the $SU(2)$ link cosine
under the Poisson-isotropic link measure contains it with locked coefficient,
\begin{equation}
E^{(4)}=-\frac{3a^4}{5760}\,S+\frac{a^4}{2880}\,K,\qquad
K=\sum_{\mu\nu}|c_\mu\times c_\nu|^2,
\end{equation}
with $S=(\mathrm{Tr}\,G)^2$, $G_{\mu\nu}=c_\mu\!\cdot\!c_\nu$; a cubic lattice is exactly
blind to $K$ (this is why lattice Skyrme terms are added by hand), and the same isotropy
that gives Lorentz invariance generates the operator [measured, with exact identities].
What is external is its \emph{dominance}---the positive net quartic that Derrick
stability needs---and the necessity proof now has two independent halves:

\emph{Tree level} [theorem]: $K\le(1-\tfrac1d)S$ is a pointwise Cauchy--Schwarz identity
on the PSD current Gram, saturated by the hedgehog; the net quartic is
$-(a^4/384)\langle|\ell|^4\rangle\le0$ under \emph{any} link measure, because its sign
descends from the cosine series, not from geometry. No regime exists to search.

\emph{One loop} [measured, $\sim$130$\sigma$]: integrating the Gaussian fluctuations
around constant-current backgrounds on the causal substrate (spatial-torus sprinkling,
per-link transporters realizing the two SC channels exactly; the abelian channel has
$K=0$, the hedgehog-frame channel $K=6g^4$ at matched $S=9g^4$) yields the radiative
channel coefficients
\begin{align}
c_K^{\rm loop}&=-0.01296\pm0.00010,\nonumber\\
c_S^{\rm loop}&=-0.00669\pm0.00007,
\end{align}
volume-stable across sizes. The loop generates the Skyrme channel with the
\emph{anti-stabilizing} sign, and its hedgehog quartic is itself negative; since the
loop term is coupling-independent while the (negative) tree term scales with $J$, the
net hedgehog quartic is negative \emph{for every} $J\ge0$---there is no window at any
coupling, ordered or not. The loop is in fact more hedgehog-saturating than the tree
(channel ratio $\approx2.3$ against the tree's $5/9$), and the contrast inverts on the
grid: the cubic control gives a weakly \emph{positive} $c_K^{\rm loop}$
($+0.001$)---the same isotropy that generates the operator at tree level kills its
dominance radiatively.

The stabilizer (a non-cosine core cost) is therefore a genuine axiom-level input of the
matter sector, located by impossibility arguments at both orders rather than by failed
searches.

\subsection{The single unit}

The second external input is the conversion of the lattice unit to SI. The thinning RG
(Sec.~\ref{sec:thinning}) shows $\rho$ is the unique marginal direction and pure gauge
at the level of the intrinsic law: the theory has exactly one unit, so exactly one
number must be imported---as in any physical theory---while every internal ratio is
locked (measured examples across the programme: the Skyrme length $a/\sqrt{120}$, the
emergent $G\propto1/K$, $m_A\propto\sqrt\rho$). The structural point is not that the
input is avoidable but that there is provably only \emph{one} of it.

\section{Where the principle does not deliver}
\label{sec:vaccine}

An exclusion principle earns trust by failing to explain things it should not explain.

\textbf{Horizon entropy, two faces [measured].} The geometric face obeys the principle's
pattern: the horizon-molecule count reproduces $S\propto A$ in $d=2,3,4$ with the
correct $\rho^{1/2}$ scaling and $T$-independence---the \emph{form} emerges; the
coefficient (the $1/4$) is a lattice number, external, exactly as the scale section
predicts. The matter face does \emph{not}: the mutual information of the ordered
internal field across the same horizon is super-area ($I\sim L^{3.4}$, compatible with
volume), driven by the layer-1 non-locality (boundary-crossing links grow $\sim L^{4.2}$).
The binary tension that forbids the photon also breaks the area law for matter
entropy---the principle predicts its own failure mode here, but it \emph{is} a failure
of the area law, and we display it rather than absorb it. The covariant quantization
refines this face without overturning it: the spacelike entanglement entropy of the SJ
field across the same horizon, computed with a declared UV number-cutoff, \emph{localizes}
on the boundary---the truncated entropy is window-independent (ratio $0.93$ against $4.0$
for the untruncated volume law), recovering the \emph{qualitative} content of an area law
that the classical mutual information had lost. But the quantitative scale is set by the
cutoff, not the geometry: the exponent tracks the truncation number $n\propto N_O^{3/4}$
and stays sub-area ($\kappa=1.57\pm0.02$), while the untruncated law reproduces the
classical super-volume [measured]. Approach to a genuine area law from below is not
excluded (the curvature rises with region size), and the discriminating grid
$L\in\{4,6\}$ is a declared open frontier. We record the refinement without upgrading it:
between a full quantum rescue of the area law and the raw matter super-area, the
measurement sits in the middle.

\textbf{Dynamics.} Ordering temperatures, phase diagrams, the first-order character of
the $SU(3)$ transition: measured, not derived. The class fixes permissions and forms.

\textbf{Quantum structure.} The substrate is classical-statistical, and the only
quantization compatible with the class is covariant---canonical quantization has no
carrier (P1), so the Sorkin--Johnston (SJ) prescription is forced. That door is now
measured to open: on the programme's box substrate the SJ vacuum is stable (conjugate
eigenvalue pairs to $10^{-15}$) and its infrared branch is relativistic, $\omega=c|k|$
with $c=0.96$ ($d=2$) and $0.90$ ($d=4$), a massive control tracking to $\sim\!2\%$
[measured, pre-registered]. The structural cost is exactly \emph{one bit}: the
retarded/advanced choice conjugates the vacuum anti-unitarily (a Wigner $T$-type
symmetry), so every spectral and entropy observable is bit-even and the campaign never
sweeps it [theorem]---the same bit that carries the complex structure the GOE corollary
(Sec.~\ref{sec:predictions}) sees as \emph{not} emergent. What the core still says
nothing about is the \emph{value} of $\hbar$, collapse, and measurement.

\section{Predictions}
\label{sec:predictions}

The core generates consequences beyond the four known negatives; each is falsifiable
inside or outside the programme.

\begin{enumerate}
\item[P1.] \emph{Canonical quantization has no carrier; covariant quantization is
forced.} Everything defined on spacelike surfaces falls at once: ADM variables,
Wheeler--DeWitt wavefunctionals, Hamiltonian lattice gauge theory, spin networks as
fundamental kinematics, spatial Wilson lines. What remains is histories/covariant
(Sorkin--Johnston, decoherence functionals). \emph{Now partly measured}: the forced SJ
quantization has a stable vacuum with a relativistic infrared branch on the programme's
substrate, at the cost of one $T$-bit (Sec.~\ref{sec:vaccine}) [measured].
\item[P2.] \emph{No $P$ violation in pure-geometry observables, ever}---including
spontaneous. A chiral gravitational-wave background of substrate origin, or $P$-odd
primordial correlators of the geometric sector, are impossible in the class.
\item[P3.] \emph{No dynamical baryogenesis.} $B$ is exactly conserved ($\pi_3$); creation
channels measured dead; the $P$-entangled content of $CP$ is external. Matter--antimatter
asymmetry is initial condition or input.
\item[P4.] \emph{No magnetic monopoles, twice}: no emergent $U(1)$, and $\pi_2(G)=0$ for
every Lie group (cosets with $\pi_2\neq0$ would need an unbroken $U(1)$---the same closed
door).
\item[P5.] \emph{No emergent stable particle without a discrete topological index; no
continuous emergent charge} (the Wigner pattern, layer 3).
\item[P6.] \emph{Hierarchy of possible spacetime symmetry breakings}: only sectors with
carriers can ever condense---$T$ (order duality) or frame/boost alignment; never $P$.
\item[P7.] \emph{Chirality, in any order-based substrate, can only be a relational
phase} (domains, $\mathbb{Z}_2$ walls, restorable), never a law. The exact, never-restored
chirality of the neutrino is structural evidence against any causal-substrate origin of
chirality.
\end{enumerate}

Two bonus corollaries of the same layers, upgraded from measurements: the causal
spectral statistics must be the \emph{orthogonal} (GOE) class---$T$-statistical symmetry
plus no complex structure (no photon) force real symmetric operators [sketch; measured
$\langle r\rangle\to0.53$]; and any unsigned counting observable must be an \emph{even}
function of an applied field-like parameter, fixing the measured $H^2$ leading behaviour
of interval anisotropies [sketch; measured $R^2>0.98$].

\section{The open frontiers}
\label{sec:open}

\begin{table}[t]
\caption{\label{tab:open}The honest list. ``Open'' means: declared, named, and either
queued or blocked on a stated obstacle.}
\begin{ruledtabular}
\begin{tabular}{clp{4.2cm}}
\# & Frontier & Status \\
\colrule
1 & pentagonal loop lattices & the single combinatorial flank (rank-frustrated
generators); named target \\
2 & Axiom 2 (internal field) & postulated; measured that the SJ vacuum creates no
compact fiber---delimitation, not derivation \\
3 & Skyrme stabilizer & external; necessity proved at tree (theorem) and one loop
(measured) \\
4 & quantization / $\hbar$ & SJ door measured open (stable vacuum, relativistic branch,
one $T$-bit); value of $\hbar$, collapse open \\
5 & entropy coefficient $1/4$ & external (lattice numbers measured) \\
6 & joint substrate+matter RG & analytic EFT; old open item \\
7 & CPT of the emergent EFT & assumed from continuum \\
8 & arrow of time ($T$ side) & closed empirically, not kinematically \\
9 & dynamical geometry & outside Axioms 1--2 (Lemma~\ref{lem:quenched} does not
cover it) \\
10 & gravitational dynamics & not derived (the emergent Newtonian sector is what
exists) \\
\end{tabular}
\end{ruledtabular}
\end{table}

Table~\ref{tab:open} is the boundary of the theory as of this writing. Item 1 is the
only place where the combinatorial theorems of Sec.~\ref{sec:trichotomy} leave a
constructive escape; items 2 and 4 are where new physics would enter if it enters;
items 5--7 are the technical debts; items 8--10 are the declared scope walls.

\section{How to break this theory}
\label{sec:conclusion}

The core is two axioms, one pair lemma, one three-layer boundary theorem, a
combinatorial trichotomy, a thinning RG, a parity no-go, and two proven-external inputs.
Everything else in the programme is either a corollary, a measurement, or listed in
Table~\ref{tab:open}. Because the core is deductive, it is attackable at named joints:
an emergent photon on an invariant causal substrate must break
Theorem~\ref{thm:T1}--\ref{thm:T4} through the pentagonal flank or falsify
Lemma~\ref{lem:pair}; emergent chirality must break Lemma~\ref{lem:isometries} or
\ref{lem:gram}; an emergent scale must exhibit a relevant direction under thinning,
against Sec.~\ref{sec:thinning}; a radiatively stabilized Skyrmion must overturn a
$130\sigma$ measurement or find physics beyond one loop with an inverted sign. A theory
that tells you exactly where to hit it is doing its job; this one, at least, now fits on
a page.

\appendix

\section{Proofs of the base lemmas}
\label{app:proofs}

\emph{Pair lemma (timelike half).} The orthochronous Poincar\'e group acts transitively
on ordered timelike pairs of equal proper time (translate $x$ to the origin, boost
$y-x$ to $(\Delta\tau,\vec0)$, rotate), with stabilizer $SO(d-1)$: the complete
invariant of an ordered causal pair is $\Delta\tau$, and the order itself is the sign.
The Alexandrov interval has volume $c_d\Delta\tau^d$, so
$N\sim\mathrm{Poisson}(\rho c_d\Delta\tau^d)$; the longest-chain limit exists by
Kingman's subadditive ergodic theorem, with constant $m_d$
bounded in~\cite{BrightwellGregory1991} ($m_2=2$ exactly, via Ulam's problem).

\emph{Pair lemma (spacelike half).} Place the pair at $0$ and $(0,s,\vec0)$. The
stabilizer of the spacelike separation vector is the orthogonal group of its Lorentzian
complement and contains the transverse reflections
$(t,x^1,\vec u)\mapsto(t,x^1,R\vec u)$, $R\in O(d-2)$, $\det R=-1$: time-orientation
preserving, spatial determinant $-1$. Each is an order isomorphism of every realization
and preserves the Poisson law, so every intrinsic statistic of the pair-plus-environment
is invariant under it---no transverse orientation is intrinsic. Distance
reconstructions~\cite{RideoutWallden2009} are symmetric functions of counts, and a
configuration known through mutual distances is fixed only up to the full orthogonal
group (Gram rigidity), reflections included.

\emph{Exhaustion (Lemma~\ref{lem:gram}).} Let $O$ be a measurable function of the
isomorphism class of $(C,\prec,n)$ and $g$ any time-orientation-preserving isometry,
spatial reflections included. Then $g$ preserves cones and time orientation, so
$x\prec y\Leftrightarrow gx\prec gy$: $g$ realizes an isomorphism
$(C_\Phi,\prec,n)\cong(C_{g\Phi},\prec,n\circ g^{-1})$, whence $O(g\Phi)=O(\Phi)$
realization-wise. The induced action of $P$ on intrinsic data is the identity; a
$P$-odd $O$ satisfies $O=-O\equiv0$. $\square$

\emph{No chiral LRO (Lemma~\ref{lem:quenched}).} Let $A,B$ be disjoint bounded regions
and $\chi_A$ a functional of $\Phi\cap A$ odd under a Lebesgue-preserving,
time-preserving spatial reflection $r_A$ of $A$. Complete independence of the Poisson
process gives $\mathbb{E}[\chi_A\chi_B]=\mathbb{E}[\chi_A]\mathbb{E}[\chi_B]$;
invariance of the restricted law under $r_A$ gives
$\mathbb{E}[\chi_A]=\mathbb{E}[\chi_A(r_A\Phi_A)]=-\mathbb{E}[\chi_A]=0$. $\square$

\emph{$\delta_0$ lemma.} The orbits of $SO^{+}(d-1,1)$ on $\mathbb{M}^d$ are $\{0\}$,
the mass shells, the light-cone components, and the one-sheeted hyperboloids; each
non-trivial orbit is $G/H$ with $G$ semisimple (unimodular) and $H$ unimodular, so by
Weil's uniqueness theorem~\cite{Folland1995} it carries a $G$-invariant measure unique
up to scale---the standard Lorentz-invariant measure, of infinite total mass. No
non-trivial orbit supports an invariant probability; disintegrating any invariant
probability over the orbit space puts all mass on $\{0\}$: $\mu=\delta_0$. $\square$

\emph{Thinning fixed point and eigenvalues.} On Laplace functionals, thinning acts as
$L\mapsto L[-\log(1-p(1-e^{-f}))]$; for $\mathrm{PPP}(\rho)$,
$L[f]=\exp(-\rho\!\int(1-e^{-f}))$, thinning gives intensity $p\rho$ and the
compensating contraction $S_{p^{1/d}}$ restores it: $T_p\,\mathrm{PPP}(\rho)=
\mathrm{PPP}(\rho)$ exactly. On factorial moment densities, thinning multiplies
$\rho_k$ by $p^k$ and $S_a$ maps $\rho_k(x)\mapsto a^{-kd}\rho_k(x/a)$; with $a^d=p$
the amplitudes cancel and $\rho_k\mapsto\rho_k(\cdot/a)$---pure argument dilation.
Homogeneous truncated directions $g=|x|^{-\alpha}$ obey $T_p\,g=p^{\alpha/d}g$.
$\square$

\begin{thebibliography}{99}
\bibitem{MatterGravityPRD} M. Alves Mendes, \emph{Topological matter and emergent
gravitation from a causal network}, companion paper, submitted to Phys. Rev. D (2026);
source at \url{https://github.com/mendesengproj-blip/TEIC}.
\bibitem{SU3PRD} M. Alves Mendes, \emph{Colour ferromagnetism in a causal network:
confinement, the meson octet, and the order of the transition}, companion paper,
submitted to Phys. Rev. D (2026); source at
\url{https://github.com/mendesengproj-blip/TEIC}.
\bibitem{WenComplement} M. Alves Mendes, \emph{Why Lorentz-invariant causal sets cannot
support an emergent $U(1)$: a causal complement to string-net theory}, companion paper,
submitted to Class. Quantum Grav. (2026); source at
\url{https://github.com/mendesengproj-blip/causal-substrate-core}.
\bibitem{BHS2009} L. Bombelli, J. Henson, R. D. Sorkin, \emph{Discreteness without
symmetry breaking: a theorem}, Mod. Phys. Lett. A \textbf{24}, 2579 (2009).
\bibitem{Malament1977} D. B. Malament, \emph{The class of continuous timelike curves
determines the topology of spacetime}, J. Math. Phys. \textbf{18}, 1399 (1977).
\bibitem{HKM1976} S. W. Hawking, A. R. King, P. J. McCarthy, \emph{A new topology for
curved space-time which incorporates the causal, differential, and conformal
structures}, J. Math. Phys. \textbf{17}, 174 (1976).
\bibitem{Surya2019} S. Surya, \emph{The causal set approach to quantum gravity},
Living Rev. Relativ. \textbf{22}, 5 (2019).
\bibitem{LevinWen2005} M. A. Levin, X.-G. Wen, \emph{String-net condensation: A physical
mechanism for topological phases}, Phys. Rev. B \textbf{71}, 045110 (2005).
\bibitem{Wigner1939} E. Wigner, \emph{On unitary representations of the inhomogeneous
Lorentz group}, Ann. Math. \textbf{40}, 149 (1939).
\bibitem{Janson2011} S. Janson, \emph{Poset limits and exchangeable random posets},
Combinatorica \textbf{31}, 529 (2011).
\bibitem{Brightwell1993} G. Brightwell, \emph{Models of random partial orders}, in
\emph{Surveys in Combinatorics 1993}, Cambridge Univ. Press (1993).
\bibitem{RideoutSorkin2000} D. P. Rideout, R. D. Sorkin, \emph{Classical sequential
growth dynamics for causal sets}, Phys. Rev. D \textbf{61}, 024002 (2000).
\bibitem{ABBJ1994} N. Alon, B. Bollob\'as, G. Brightwell, S. Janson, \emph{Linear
extensions of a random partial order}, Ann. Appl. Probab. \textbf{4}, 108 (1994).
\bibitem{BollobasBrightwell1997} B. Bollob\'as, G. Brightwell, \emph{The structure of
random graph orders}, SIAM J. Discrete Math. \textbf{10}, 318 (1997).
\bibitem{Kallenberg2017} O. Kallenberg, \emph{Random Measures, Theory and Applications},
Springer (2017).
\bibitem{Elitzur1975} S. Elitzur, \emph{Impossibility of spontaneously breaking local
symmetries}, Phys. Rev. D \textbf{12}, 3978 (1975).
\bibitem{BrightwellGregory1991} G. Brightwell, R. Gregory, \emph{Structure of random
discrete spacetime}, Phys. Rev. Lett. \textbf{66}, 260 (1991).
\bibitem{RideoutWallden2009} D. Rideout, P. Wallden, \emph{Spacelike distance from
discrete causal order}, Class. Quantum Grav. \textbf{26}, 155013 (2009).
\bibitem{Folland1995} G. B. Folland, \emph{A Course in Abstract Harmonic Analysis},
CRC Press (1995), Thm. 2.49.
\end{thebibliography}

\end{document}
